\section{Introduction}
Chart summarization requires faithfully extracting quantitative data and describing them using natural language. Recent natural language generation (NLG) studies have explored different flavors of chart-to-summary generation tasks including caption generation for scientific figures \cite{hsu-etal-2021-scicap-generating}, chart summary generation \cite{kantharaj-etal-2022-chart}, or analytical textual descriptions for charts \cite{zhu-etal-2021-autochart}. These tasks can be advantageous for the visually impaired \citep{benetech-making-graphs-accessible} as well as for automating interpreting complex domains such as finance data-analysis, news reporting, and scientific domains \citep{siegel2016figureseer}.

\begin{figure}[!t]
    \centering
    \includegraphics[width=1.\linewidth]{figures/ChaTS-Pi.pdf}
    \caption{\pipeline generates multiple summaries given the chart using any summarization model. Each summary is then \emph{repaired} by dropping refuted sentences according to the \metric{} sentence scoring. 
    Finally, we rank the summaries by computing the ratio of sentences that were kept.
    }
    \label{fig:ChaTS-Pi}
    \vspace{-2.0ex}
\end{figure}

While a wide range of models and techniques have been applied for chart summarization, faithfulness remains a major challenge for the task.
Specifically, the models often misread details in the charts (due to perceptual mistakes) or miscalculate the aggregations (due to reasoning flaws).
To overcome some of these limitations, \emph{optical character recognition} (OCR) models and object detection systems are usually employed to extract meta-data such as axis, values, titles, legend~\cite{luo2021chartocr,masry-etal-2022-chartqa}. These data are then used as auxiliary inputs to finetune NLG models. Nonetheless, these modeling efforts are still limited by two fundamental issues (i) training \& evaluation dataset quality and (ii) the reference-based metrics being used for evaluation.
As examples, two widely used datasets, Chart-to-Text~\cite{kantharaj-etal-2022-chart} and SciCap~\cite{hsu-etal-2021-scicap-generating}, are automatically extracted from web articles and academic journals. As a result, the summary references are prone to \emph{hallucination}, i.e. the reference might contain context that cannot be entailed solely by the chart content. Training on this data can encourage the NLG models to improvise/hallucinate. Besides, the auto-extracted summaries sometimes emphasize only certain aspects of the chart, missing out critical insights.

On the other hand, n-gram based metrics such as BLEU~\cite{papineni-etal-2002-bleu}, or learned metrics such as BLEURT~\cite{sellam-etal-2020-bleurt} rely only on gold references. They are not capable of recognizing 
unreferenced but correct insights since they solely rely on the reference for scoring the summaries, 
as shown in \Cref{fig:example_statista}. This issue is especially pronounced when the gold references are noisy, which is the case for Chart-to-Text and SciCap.
Last but not least, reference-based metrics also heavily penalize summary style mismatches, giving an artificial disadvantage to LLMs which are not tuned on the task data ~\cite{maynez-etal-2023-benchmarking}.

This motivates building a reference-free critic \metric{} (\Cref{fig:CHATS-CRITIC}) that can be used as a metric
to score and re-rank summaries.
We additionally introduce \pipeline{} (\Cref{fig:ChaTS-Pi}) that leverage \metric{} scores to generate a high quality summaries.
We summarize our contributions as follows:

\begin{figure}[!ht]
 \vspace{-1.0ex}
    \centering
    \includegraphics[width=0.90\linewidth]{figures/CHATS-CRITIC.pdf}
  \vspace{-2.0ex}
    \caption{\metric{} is composed of a de-rendering model to extract the table from the chart, and a table entailment model. The latter can be a blackbox table entailment model (e.g., TabFact as benchmarked in \Cref{tab:chats-critic-size}) or an LLM; in latter case, we use CoT prompt and average over~$8$ samples. In the figure, the threshold to reach a binary decision is set to $T=0.75$. The chart icon refers to the same plot of \Cref{fig:example_statista}.}
    \label{fig:CHATS-CRITIC}
    \vspace{-2.0ex}
\end{figure}

\begin{enumerate}[noitemsep]
    \item We present \metric{}, a reference-free metric composed of a model that extracts the underlying table data from the chart and a table-entailment model acting on each sentence within a chart summary.
    \item We design \pipeline{}, a pipeline that (i) generates multiple candidate summaries using a generative model, either fine-tuned or with in-context learning; (ii) then leverages \metric{} to refine the summaries by dropping unsupported sentences; (iii) computes a summary score to rank the summaries by penalizing summaries with dropped sentences to increase the fluency, and (iv) outputs the best one.
    \item To assess the efficacy of \metric{}, we juxtapose human preferences against both \metric{} and other prevailing metrics. Our results indicate that \metric{} aligns more consistently with human evaluations. Furthermore, when contrasting \pipeline{} with other leading models that serve as baselines, \pipeline{} establishes \sota on two popular English benchmarks.

\end{enumerate}


\begin{figure}[!ht]
    \centering
    \includegraphics[width=1.0\linewidth]{figures/example_limits_current_models.pdf}
    \vspace{-4ex}
    \caption{This example from \citet{kantharaj-etal-2022-chart} showcases the limits of reference-based metrics for summary evaluation: (1) the reference text often contains extra information that is not present in the chart which skews the evaluation, and (2) the reference-based metrics can fail at capturing unreferenced but correct sentences. In comparison, \metric{} better reflects the human ratings for summary faithfulness.
    }
    \label{fig:example_statista}
    \vspace{-3ex}
\end{figure}
